\subsubsection{Péndulo en collarín móvil}

Una partícula de masa $m$ está suspendida desde el punto $P$, ubicado en un extremo de una varilla esbelta, rígida y sin masa, como se muestra en la **Figura \ref{fig:pendulo_collarin}**. El otro extremo de la varilla está articulado en un punto $Q$, donde $Q$ está unido a un collarín sin masa. El collarín se desliza con un desplazamiento especificado $x(t)$ a lo largo de una guía horizontal. Sabiendo que $\theta$ describe la orientación de la varilla con respecto a la dirección vertical hacia abajo y que la gravedad actúa hacia abajo.

\begin{enumerate}
    \item Determine la ecuación diferencial de movimiento de la partícula utilizando la segunda ley de Newton.
    \item Determine la ecuación diferencia de movimiento de la partícula utilizando momento angular.
\end{enumerate}

\begin{figure}[!ht]
    \centering
    \begin{tikzpicture}[>=Latex, scale=1.3]
        % Guía horizontal
        \fill[pattern=north east lines] (-3, 0.2) rectangle (3, 0.3);
        \fill[pattern=north east lines] (-3, -0.3) rectangle (3, -0.2);
        \draw[thick] (-3, 0.2) -- (3, 0.2);
        \draw[thick] (-3, -0.2) -- (3, -0.2);
        \fill[top color=gray!40, middle color=white, bottom color=gray!40] (-3, -0.2) rectangle (3, 0.2);

        % Collarín Q
        \shadedraw[ball color=white!90, draw=black] (0, -0.15) rectangle (0.8, 0.15);
        \coordinate (Q) at (0.4, 0);
        \draw[dashed] (Q) -- (0.4, -1.5);
        \filldraw (Q) circle (2pt) node[below right] {$Q$};
        \draw[->, thick] (-1.4, 0.6) -- (0.4, 0.6) node[midway, above] {$x(t)$};
        \draw[gray] (-1.4, 0.3) -- (-1.4, 0.7);

        % Varilla y Masa
        \begin{scope}[shift={(Q)}, rotate=-35]
            \draw[ultra thick] (0,0) -- (0, -2.5) node[midway, right] {$l$};
            \draw[thick, color=gray] (0,0) -- (0, -2.5);
            \shadedraw[ball color=white] (0, -2.5) circle (5pt) 
                node[right, xshift=4pt] {$P$} 
                node[font=\footnotesize] {$m$}; % Tamaño ajustado aquí
        \end{scope}

        % Ángulo theta
        \draw[->] (0.4, -0.8) arc (-90:-125:0.8);
        \node at (0.2, -0.6) {$\theta$};
        
        % Gravedad
        \draw[->, thick] (-1.5, -0.8) -- (-1.5, -1.8) node[midway, left] {$\mathbf{g}$};
    \end{tikzpicture}
    \caption{Esquema cinemático de un péndulo de masa $m$ y longitud $l$ acoplado a un collarín con movimiento prescrito $x(t)$.}
    \label{fig:pendulo_collarin}
\end{figure}